\section{Scope conditions: when the frame degrades}
\label{sec:scope}

\DCI\ is not a universal solvent, and its guarantees are only as strong as its two assets: a compressed vertically organizing principle $\Prin$ and an executable verifier $\Rule$. Both can be partial, and the frame should --- and does --- degrade gracefully and measurably as either weakens. Stating where it fails is part of what makes it falsifiable.

\emph{Weak or absent $\Prin$.} Many domains of interest may not admit a single scale-invariant action functional; much of biology, economics, and social dynamics may offer only regional symmetries, approximate conservation laws, or learned invariants rather than a genuine $\Prin$. There, \emph{reflexivity through inference} (\S\ref{sec:cross}, Instance~2) weakens with it: absent a true $\Prin$, the meta-level objective is no longer guaranteed to share the object-level's variational form, and the object/meta boundary no longer collapses. The frame makes no claim of domain-completeness in this regime; it predicts \emph{partial} completeness in proportion to how much lawful structure $\Prin$ compresses. This is exactly prediction~(ii) of \S\ref{sec:testable}: a domain with no compressible invariant should show no \DCI\ advantage, which is a way to be wrong, not a hedge.

\emph{Weak or approximate $\Rule$.} Where the substrate is only partly executable --- noisy simulators, costly real-world experiments, oracles that return a signal rather than ground truth --- the closed loop leaks: false positives survive verification and the cost of experience $E$ stops being cheap. Here the composition of the two assets (\S\ref{sec:cross}) does the work: $\Prin$'s meta-coverage must bear more weight precisely where $\Rule$ is unreliable, and $\Rule$ must carry more where $\Prin$ is thin. When \emph{both} are weak at once, \DCI\ has no special leverage and collapses toward ordinary narrow AI, or toward the open-world agentic problem it was meant to precede. That boundary is the honest edge of the proposal, not a concealed failure of it.

\emph{Quantifying the leak.} ``Degrades gracefully'' is only honest if the degradation has a stated shape, so we give one. Let the verifier have fidelity $\rho$ --- the probability it adjudicates a proposal correctly, equivalently one minus its false-positive rate on plausible-but-wrong candidates --- and a per-query cost $c$: a perfect executable substrate has $\rho = 1,\ c \to 0$, while a stochastic simulator or a low-throughput physical assay has $\rho < 1$ and $c$ large enough to inflate the effective experience term $E$. Two regimes bound the behaviour. As $\rho \to 1$ and $c \to 0$ the loop is the closed, cheap-$E$ discovery loop of \S\ref{sec:definition}. As $\rho$ falls, wrong laws survive verification and accumulate; below some domain-dependent threshold $\rho^\ast$ their survival rate exceeds the prior's power to suppress them, the adjudication signal is no longer trustworthy, and \DCI\ reduces to standard RL against an expensive, unreliable reward --- the closed loop provides no advantage over model-free trial-and-error. We do not have $\rho^\ast$ in closed form; locating it, and in particular its dependence on the compression of $\Prin$ (a stronger prior should tolerate a weaker verifier, since it needs fewer bits of ground truth to discriminate candidates), is a concrete measurement this program owes and a further way for the frame to be wrong. The physics case gives the intuition and a lower anchor: wide-stencil preprocessing had to lift an $\mathrm{SNR}\approx 0.02$ substrate to ${\approx}\,1.6$ \citep{frasch2026b} \emph{before} the discovery loop closed at all --- direct evidence that a verifier-quality floor exists and that discovery does not occur beneath it.

\emph{The residual burden on reflexivity.} Even granting $\Prin$ and $\Rule$, the reflexivity claim is best-evidenced for search (NAS) and for the variational form of inference; it is weakest for open-ended meta-skills such as designing a genuinely novel experimental apparatus, or recognizing that the domain's rules themselves require revision --- the basis-expansion step flagged in \S\ref{sec:verifier}. We do not claim these are fully covered by $\Prin$. We claim two narrower things: that the executable verifier substitutes for meta-judgment in the cases that matter for discovery, and that whatever remains uncovered is precisely the frontier where \DCI\ shades into the universal-agency problem --- a frontier we are content to leave to that program rather than to overclaim for this one.
